\documentclass[conference]{IEEEtran}
\usepackage[utf8]{inputenc}
\usepackage[T1]{fontenc}
\usepackage{cite}
\usepackage{url}
\usepackage{graphicx}
\usepackage{hyperref}
\usepackage{listings}
\usepackage{xcolor}

\hypersetup{colorlinks=true, urlcolor=blue, linkcolor=blue, citecolor=blue}
\sloppy

\lstset{
    basicstyle=\ttfamily\scriptsize,
    breaklines=true,
    frame=single,
    columns=fullflexible,
    backgroundcolor=\color{gray!10}
}

\begin{document}

\title{Fragmentaci\'on Horizontal y Replicaci\'on en Bases de Datos Distribuidas: \\
Implementaci\'on del Caso Banco del Austro}

\author{
\IEEEauthorblockN{Harold Nicolas Vinueza S\'anchez}
\IEEEauthorblockA{Ingenier\'ia en Software, Universidad T\'ecnica Estatal de Quevedo\\
Aplicaciones Distribuidas --- S\'eptimo Nivel, Paralelo A}
}

\maketitle

\begin{abstract}
Este trabajo documenta el dise\~no e implementaci\'on de un prototipo did\'actico de base
de datos distribuida basado en el caso Banco del Austro, entidad financiera con
sedes en Cuenca y Quito. Se implementa fragmentaci\'on horizontal sobre las tablas
\texttt{cuentas} y \texttt{transacciones}, y replicaci\'on completa de la tabla
\texttt{clientes}, sobre dos motores PostgreSQL 16 independientes coordinados por
una aplicaci\'on Spring Boot que enruta cada consulta seg\'un el prefijo del n\'umero
de cuenta. Se verifican emp\'iricamente las propiedades de transparencia de
fragmentaci\'on, transparencia de ubicaci\'on y autonom\'ia local, esta \'ultima
mediante una prueba controlada de ca\'ida de nodo.
\end{abstract}

\begin{IEEEkeywords}
bases de datos distribuidas, fragmentaci\'on horizontal, replicaci\'on, autonom\'ia
local, transparencia de ubicaci\'on, PostgreSQL, Spring Boot.
\end{IEEEkeywords}

\section{Introducci\'on}

Las organizaciones con presencia en m\'ultiples sedes f\'isicas enfrentan un
compromiso entre la necesidad de que cada oficina opere con baja latencia sobre
sus datos locales y la necesidad de que la informaci\'on cr\'itica --- como la
identidad de los clientes --- sea visible de forma consistente en toda la
organizaci\'on \cite{kleppmann2017}. El caso del Banco del Austro, con oficinas
en Cuenca, Quito, Guayaquil, Loja y Machala, ilustra este compromiso a escala
reducida: se implementa un prototipo con dos sedes (Cuenca y Quito) que preserva
las tres propiedades centrales de una base de datos distribuida --- fragmentaci\'on
horizontal, transparencia de ubicaci\'on y autonom\'ia local --- descritas
formalmente por \"Ozsu y Valduriez \cite{ozsu2020}.

El objetivo de este documento es describir el dise\~no de fragmentaci\'on y
replicaci\'on adoptado, la arquitectura de la implementaci\'on sobre contenedores
Docker y una aplicaci\'on Spring Boot con m\'ultiples \texttt{DataSource}, y
presentar la verificaci\'on emp\'irica de las propiedades distribuidas mencionadas,
incluyendo una prueba de tolerancia a fallos.

\section{Marco Conceptual}

Siguiendo la terminolog\'ia de \"Ozsu y Valduriez \cite{ozsu2020}, una \emph{sede}
es una ubicaci\'on f\'isica donde reside un motor de base de datos con parte o
la totalidad del dato de la organizaci\'on; un \emph{nodo} es cada motor
individual considerado como participante del sistema distribuido. Un
\emph{fragmento horizontal} es un subconjunto de filas de una tabla l\'ogica
completa que reside en un nodo espec\'ifico, y debe cumplir tres propiedades:
completitud (ninguna fila se pierde), reconstructibilidad (la uni\'on de los
fragmentos recupera la tabla original) y disjunci\'on (ninguna fila reside en
dos fragmentos simult\'aneamente). Una \emph{r\'eplica}, en cambio, es una copia
completa de una tabla mantenida en m\'as de un nodo para acelerar lecturas o
tolerar la ca\'ida de un nodo, a costa de encarecer las escrituras al requerir
actualizar todas las copias \cite{kleppmann2017}.

\section{Dise\~no de Fragmentaci\'on y Replicaci\'on}

El prototipo reduce el problema real del banco a tres tablas: \texttt{clientes},
\texttt{cuentas} y \texttt{transacciones}. La Tabla \ref{tab:estrategia} resume
la estrategia de distribuci\'on adoptada para cada una.

\begin{table}[h]
\centering
\caption{Estrategia de distribuci\'on por tabla}
\label{tab:estrategia}
\begin{tabular}{|p{1.6cm}|p{1.6cm}|p{3.5cm}|}
\hline
\textbf{Tabla} & \textbf{Estrategia} & \textbf{Justificaci\'on} \\
\hline
clientes & Replicada & Cualquier oficina debe poder identificar al titular de
una cuenta, incluso si pertenece a otra sede. \\
\hline
cuentas & Fragmentada horizontalmente & Cada cuenta pertenece a exactamente
una oficina; se consulta casi siempre localmente. \\
\hline
transacciones & Fragmentada horizontalmente & Cada movimiento se registra donde
se origin\'o; los movimientos entre oficinas se difieren a trabajo futuro
(commit en dos fases). \\
\hline
\end{tabular}
\end{table}

Para que la aplicaci\'on determine la sede de una cuenta sin consultar la base
de datos, se adopt\'o una convenci\'on de prefijos sobre el n\'umero de cuenta:
las cuentas de Cuenca inician con \texttt{22} y las de Quito con \texttt{17},
tomando como referencia los d\'igitos provinciales de la c\'edula ecuatoriana.
La regla de disjunci\'on de la fragmentaci\'on se garantiza a nivel de motor
mediante una restricci\'on \texttt{CHECK} en cada nodo:

\begin{lstlisting}[language=SQL]
-- Nodo Cuenca
oficina VARCHAR(20) NOT NULL
    CHECK (oficina = 'CUENCA')

-- Nodo Quito
oficina VARCHAR(20) NOT NULL
    CHECK (oficina = 'QUITO')
\end{lstlisting}

Con esta restricci\'on, cualquier intento de insertar una fila con
\texttt{oficina = 'QUITO'} en el nodo Cuenca es rechazado autom\'aticamente
por PostgreSQL, trasladando una regla de dise\~no l\'ogico a una garant\'ia
f\'isica del motor.

\section{Arquitectura de la Implementaci\'on}

\subsection{Capa de datos: dos motores PostgreSQL en Docker}

Cada sede se implement\'o como un contenedor Docker independiente basado en la
imagen \texttt{postgres:16-alpine}, con su propio volumen persistente y su
propio esquema inicializado autom\'aticamente mediante scripts SQL montados en
\texttt{docker-entrypoint-initdb.d} \cite{docker2024}. Durante la configuraci\'on
se detect\'o un conflicto entre el puerto est\'andar de PostgreSQL (5432) y una
instalaci\'on nativa preexistente en el equipo de desarrollo; se resolvi\'o
remapeando los puertos expuestos a \texttt{5442} (Cuenca) y \texttt{5443}
(Quito), sin modificar el puerto interno del contenedor. Este ajuste ilustra
un problema pr\'actico com\'un en el despliegue de sistemas distribuidos
locales: la coexistencia de m\'ultiples instancias del mismo motor en un
mismo host.

\subsection{Capa de aplicaci\'on: Spring Boot con m\'ultiples \texttt{DataSource}}

La aplicaci\'on cliente se implement\'o en Java 21 sobre Spring Boot 3.5.16,
con dos \texttt{DataSource} independientes registrados como \emph{beans}
calificados (\texttt{@Qualifier("dsCuenca")}, \texttt{@Qualifier("dsQuito")}).
Este es el punto exacto donde la distribuci\'on se materializa en el c\'odigo
cliente: un servicio de enrutamiento (\texttt{ConsultaDistribuidaService})
inspecciona los dos primeros d\'igitos del n\'umero de cuenta recibido y
selecciona el \texttt{JdbcTemplate} correspondiente antes de ejecutar la
consulta:

\begin{lstlisting}[language=Java]
private JdbcTemplate enrutar(String numero) {
    String prefijo = numero.substring(0, 2);
    return switch (prefijo) {
        case "22" -> jdbcCuenca;
        case "17" -> jdbcQuito;
        default -> throw new IllegalArgumentException(
            "Prefijo de oficina no reconocido: " + prefijo);
    };
}
\end{lstlisting}

El controlador REST expone dos endpoints: \texttt{GET /api/banco/saldo/\{numero\}},
que enruta la consulta a un \'unico nodo, y \texttt{GET /api/banco/clientes},
que consulta ambos nodos y concatena los resultados, ilustrando la operaci\'on
de uni\'on sobre fragmentos y r\'eplicas descrita por \"Ozsu y Valduriez
\cite{ozsu2020}. Ninguno de los dos endpoints revela al cliente HTTP cu\'antos
motores existen ni d\'onde reside f\'isicamente cada fila, lo que constituye
una implementaci\'on concreta de \emph{transparencia de fragmentaci\'on} y
\emph{transparencia de ubicaci\'on}.

\section{Verificaci\'on Emp\'irica}

\subsection{Verificaci\'on de la fragmentaci\'on en el motor}

Mediante pgAdmin 4 se conect\'o de forma independiente a cada nodo
(\texttt{banco\_cuenca} en el puerto 5442, \texttt{banco\_quito} en el
puerto 5443) y se ejecut\'o una consulta de uni\'on entre \texttt{cuentas} y
\texttt{clientes} en cada servidor. El nodo Cuenca devolvi\'o exactamente las
tres cuentas con prefijo \texttt{22}; el nodo Quito devolvi\'o las tres cuentas
con prefijo \texttt{17}. La suma de ambos fragmentos (3 + 3 = 6 filas) coincide
con el total esperado de la relaci\'on l\'ogica completa, verificando
emp\'iricamente el teorema de reconstrucci\'on sobre fragmentos disjuntos
\cite{ozsu2020}.

\subsection{Verificaci\'on de la transparencia de ubicaci\'on}

Se realizaron peticiones HTTP directas sobre la aplicaci\'on en ejecuci\'on
(puerto 8080), obteniendo las siguientes respuestas reales:

\begin{lstlisting}
GET /api/banco/saldo/2201000001
{"numero":"2201000001","saldo":15000.00,
 "oficina":"CUENCA"}

GET /api/banco/saldo/1701000001
{"numero":"1701000001","saldo":8000.00,
 "oficina":"QUITO"}

GET /api/banco/clientes
[6 registros: 3 de Cuenca + 3 de Quito]
\end{lstlisting}

Las tres peticiones se resolvieron correctamente sin que el cliente HTTP
especificara en ning\'un momento cu\'al motor f\'isico deb\'ia consultarse,
confirmando que el enrutador interno de la aplicaci\'on cumple su funci\'on
de forma transparente.

\section{Prueba de Tolerancia a Fallos}

Para verificar la propiedad de autonom\'ia local se detuvo deliberadamente el
contenedor \texttt{banco-quito} mediante \texttt{docker stop} mientras la
aplicaci\'on permanec\'ia en ejecuci\'on, y se repitieron las tres peticiones
anteriores. El resultado observado fue el siguiente:

\begin{itemize}
    \item La consulta sobre la cuenta \texttt{2201000001} (nodo Cuenca) se
    resolvi\'o con normalidad, con la misma respuesta que antes de la ca\'ida.
    \item La consulta sobre la cuenta \texttt{1701000001} (nodo Quito) no
    devolvi\'o respuesta ni error inmediato: la conexi\'on qued\'o en espera
    hasta ser interrumpida manualmente, comportamiento consistente con el
    agotamiento del tiempo de espera por defecto del pool de conexiones
    HikariCP, dado que la aplicaci\'on no implementa un patr\'on de tipo
    \emph{circuit breaker} \cite{newman2021}.
    \item La consulta \texttt{/api/banco/clientes}, al depender de ambos
    nodos para la operaci\'on de uni\'on, present\'o el mismo comportamiento
    de espera prolongada que la consulta a Quito.
\end{itemize}

Al reactivar el contenedor con \texttt{docker start}, las tres peticiones
volvieron a resolverse con normalidad sin necesidad de reiniciar la
aplicaci\'on Java. Este resultado confirma emp\'iricamente que el nodo Cuenca
mantuvo disponibilidad total sobre sus propios datos durante la ca\'ida de
Quito --- es decir, \emph{autonom\'ia local} --- si bien evidencia tambi\'en
una limitaci\'on real del prototipo: sin un mecanismo de \emph{circuit
breaker}, la indisponibilidad de un nodo se propaga como una espera larga
incluso hacia operaciones que, en principio, no lo necesitan \cite{newman2021}.
Este comportamiento es consistente con una clasificaci\'on CP (consistencia
y tolerancia a particiones) seg\'un el teorema CAP \cite{brewer2000}
\cite{gilbert2002}, matizada por el modelo PACELC \cite{abadi2012}: ante la
p\'erdida de un nodo, el sistema prioriz\'o no devolver datos potencialmente
inconsistentes sobre Quito antes que degradar autom\'aticamente el servicio.

\section{Aplicabilidad al Proyecto Integrador}

El proyecto integrador SCLI (Sistema de Control de Laboratorio Inform\'atico),
compuesto por m\'odulos NestJS, una base de datos PostgreSQL compartida y un
API Gateway con autenticaci\'on JWT y control de acceso basado en roles,
podr\'ia beneficiarse de un esquema de fragmentaci\'on an\'alogo si el sistema
llegara a operar sobre m\'ultiples laboratorios f\'isicos distribuidos en
distintas sedes de la instituci\'on: los registros de uso y reservas de cada
laboratorio se fragmentar\'ian horizontalmente por sede, mientras que el
cat\'alogo de usuarios permanecer\'ia replicado para permitir autenticaci\'on
cruzada entre sedes, siguiendo el mismo principio aplicado a las tablas
\texttt{cuentas} y \texttt{clientes} en este prototipo.

\section{Conclusiones}

El prototipo implementado evidencia de forma pr\'actica las tres propiedades
centrales de una base de datos distribuida: fragmentaci\'on horizontal
completa, reconstructible y disjunta --- garantizada a nivel de motor mediante
restricciones \texttt{CHECK}--- replicaci\'on de datos compartidos, y
autonom\'ia local verificada emp\'iricamente mediante una prueba de ca\'ida de
nodo controlada. La principal limitaci\'on identificada es la ausencia de un
patr\'on de \emph{circuit breaker}, que en su estado actual provoca que la
ca\'ida de un nodo se propague como una espera prolongada hacia operaciones de
uni\'on entre fragmentos. Trabajo futuro incluye la incorporaci\'on de dicho
patr\'on, la implementaci\'on de \emph{commit} en dos fases para transferencias
entre sedes y la extensi\'on del esquema a un tercer nodo (Guayaquil).

\section*{Declaraci\'on de Uso de Inteligencia Artificial}
Se utiliz\'o el asistente de inteligencia artificial Claude (Anthropic) como
apoyo durante la implementaci\'on guiada paso a paso del c\'odigo (configuraci\'on
de Docker, clases Java de Spring Boot) y en la redacci\'on y estructuraci\'on
de este documento. Todos los resultados emp\'iricos reportados corresponden a
ejecuciones reales realizadas por el autor sobre su propio ambiente de
desarrollo.

\begin{thebibliography}{00}

\bibitem{ozsu2020} M. T. \"Ozsu and P. Valduriez, \emph{Principles of
Distributed Database Systems}, 4th ed. Cham, Switzerland: Springer
International Publishing, 2020.

\bibitem{kleppmann2017} M. Kleppmann, \emph{Designing Data-Intensive
Applications: The Big Ideas Behind Reliable, Scalable, and Maintainable
Systems}, 1st ed. Sebastopol, CA, USA: O'Reilly Media, 2017.

\bibitem{newman2021} S. Newman, \emph{Building Microservices: Designing
Fine-Grained Systems}, 2nd ed. Sebastopol, CA, USA: O'Reilly Media, 2021.

\bibitem{brewer2000} E. A. Brewer, ``Towards robust distributed systems
(abstract),'' in \emph{Proc. 19th Annu. ACM Symp. Principles of Distributed
Computing (PODC '00)}, Portland, OR, USA, 2000, p. 7, doi:
10.1145/343477.343502.

\bibitem{gilbert2002} S. Gilbert and N. Lynch, ``Brewer's conjecture and the
feasibility of consistent, available, partition-tolerant web services,''
\emph{ACM SIGACT News}, vol. 33, no. 2, pp. 51--59, Jun. 2002, doi:
10.1145/564585.564601.

\bibitem{abadi2012} D. J. Abadi, ``Consistency tradeoffs in modern
distributed database system design: CAP is only part of the story,''
\emph{Computer}, vol. 45, no. 2, pp. 37--42, Feb. 2012, doi:
10.1109/MC.2012.33.

\bibitem{docker2024} Docker, Inc., ``Docker Compose specification.''
[En l\'inea]. Disponible: \url{https://docs.docker.com/compose/}.
[Accedido: 9-jul-2026].

\bibitem{spring2024} VMware Tanzu / Broadcom, ``Spring Boot Reference
Documentation.'' [En l\'inea]. Disponible:
\url{https://spring.io/projects/spring-boot}. [Accedido: 9-jul-2026].

\bibitem{postgres2024} PostgreSQL Global Development Group, ``PostgreSQL 16
Documentation.'' [En l\'inea]. Disponible:
\url{https://www.postgresql.org/docs/16/}. [Accedido: 9-jul-2026].

\end{thebibliography}

\end{document}